\scanpage{76}
\noindent Soit donc \(U_i\) un recouvrement affine de \(X_0\), et
\[
 e\in H^{0^{(\mathrm{alg})}}(X_0,\mathscr{E}xt^1(K,\mathscr{N}))
\]
représenté par des classes d'isomorphies (algébriques) d'extensions \raisebox{0.2ex}{\parbox[b]{2.8em}{\figure{%
\begin{tikzpicture}[x=1mm,y=1mm,baseline=0pt]
\node[inner sep=0pt,anchor=south] at (3.2,3.6) {$\inserted{U'_i}$};
\draw[line width=0.35pt] (0,3.3)
 .. controls (-0.1,1.5) and (1.4,1.4) .. (2.2,1.1)
 .. controls (2.8,0.9) and (3.0,0.3) .. (3.4,0)
 .. controls (3.9,1.4) and (4.9,2.6) .. (6.1,2.8);
\end{tikzpicture}%
}}} de \(U_i\) par \(\mathscr{N}\mid U_i\), classes d'isomorphies qui se recollent :

\[
 \text{soit }\unsure{\eta|n|\gamma}_{ij}\text{ un isomorphisme}\qquad
 U'_j\mid U_i\cap U_j
 \xrightarrow{\sim}
 U'_{\unsure{i|j}}\mid U_i\cap U_j\ .
\]

L'obstruction à construire des \(\eta_{ij}\) \deleted{d}tels que \(\eta_{ij}\eta_{jk}=\eta_{ik}\) est donné par \(d_2\eta\) où \(d_2\) est une certaine flèche
\[
 d_2:H^0(X_0,\mathscr{E}xt^1(K,\mathscr{N}))\longrightarrow
 H^2(X_0,\mathscr{H}om(K,\mathscr{N}))
\]

La même construction marche dans le cas analytique, et on conclut par une nouvelle application de GAGA.

\uline{Rmq} 1 : itou pour les espaces algébriques, espaces rigides analytiques ; schémas formels

\deleted{\uline{\unsure{Rmq 2|Rmqs}}} :

\noindent\hfill Bien à toi\\
\null\hfill P. Deligne
\margin{\vspace*{4.5em}\figure{\unsure{%
\begin{tikzpicture}[x=1mm,y=1mm]
\fill (0,0) ellipse[x radius=0.23,y radius=0.21];
\fill (0.62,0.92) ellipse[x radius=0.22,y radius=0.31,rotate=-36];
\end{tikzpicture}%
}}}
% archive sheet number: 75
